Zij was de schone schijn,
hij was het meesterbrein.

Messing was ooit getrouwd
met gewiekstheid vertrouwd.

Geloof in engelen
kon zij verstrengelen

met haar list en bedrog.
Dat deden zij ook toch?

Ze was lesbienne
en had een antenne

voor gelijke vrouwen,
die ze kon vertrouwen

dit met ze te delen.
Dat waren er velen.

Onthulling bleef haar angst.
Ze bleek daarbij het bangst

dat plots haar echtgenoot
uit wraak ineens besloot

haar te openbaren
met alle gevaren

voor haar leven van dien.
Hij had haar ooit gezien.

“Angst is ons kernprobleem.
Het legt een zware claim

op ons welbevinden.
Daar kunnen we vinden

eerdere invasie,
misbruik, isolatie,

hardheid, vernedering,
afbraak of afwijzing.

Leer de angst verkennen,
op diep niveau kennen,

in ons kwetsbaar lichaam
en wees vooral achtzaam

hoe angsten gedachten
en gevoelens brachten.

Angst wordt door mij gezaaid,
als dwaze ziel verdraaid,

om de macht te krijgen
door jou dood te zwijgen.

Angst werkt hartverlammend
als het ego vlammend

zich met je lijf vermengt.
Jou van je pad afbrengt

dat je vruchteloos loopt.
Het is vrees dat ons doopt

in vastberadenheid.
Omarm bescheidenheid.

Jij bent van geen gewicht.
Volg mijn innerlijk licht

dan ben je echt veilig
en is je pad heilig

tot straks de eindtijd komt.”
Emmy raakte verstomd.

Messing sprak continu.
“Vrees niet, Ik ben met u

leest men in Jesaja.
De Bhagavad Gita,

kent de afgod Krishna
die krijger Arjuna

aanspoort te gaan strijden.
Zijn vrees te vermijden.

Dat is godslastering!
God gaf als aansporing

uit liefde ons geestdrift,
staat in het heilig schrift,

om hem te erkennen
en zo te verkennen

ons vertrouwen in hem.”
Emmy meende ad rem

in beide boodschappen
hetzelfde te snappen.

Met een klap op de wang
als berispende dwang

kreeg de Abdis gelijk
en vervolgde bloemrijk:

“Kom toch tot je zinnen.
Ga niet overwinnen,

zie angst niet als gevaar.
Zie het als Gods gebaar

voor zijn troost, steun en kracht.
Zijn oneindige macht.”

Ze kreeg haar naam Messing
van hem die ze aanhing

met heel haar essentie.
Vanuit Puducherry

voer hij als scheepsarts mee
ooit bij de VOC.

Hij zat maandenlang vast
door een gebroken mast.

Naast functionarissen
en missionarissen

sprak hij ook hun tolken.
De inheemse volken

ontmoette hij door hen.
Hij werd aldaar een fan
van een wijze goeroe
tegelijk een hindoe,

die hem onderrichtte
en zo zijn leer stichtte.

Vanuit een transcriptie
schreef hij zijn profetie.

Hij was eerst tovenaar
en speelde het zo klaar

zich als een geneesheer
te tonen menig keer.

Zo kwam hij ooit aan boord.
Men geloofde zijn woord

dat men in een gilde
hem als dokter drilde.

De kwakzalver ving bot,
toen het schip het noodlot

trof in dit India.
Men kreeg malaria

en zag zijn misleiding
toen men massaal dood ging.

Hij vluchtte via land.
Als een komediant

verdiende hij de kost
als pierrot uitgedost.

Hij koppelde humor
aan pieuze horror.

Messing trof hem pas aan,
toen hij zich voort liet gaan

als een nieuwe profeet.
Ze voelde naast angstzweet
ook een aantrekkingskracht.
Hij kreeg haar in zijn macht

door zijn trickster aura.
Zijn clownesk charisma

met een te brede grijns.
die quasi stoïcijns

zich loszinnig gedroeg.
Waarbij zij zich afvroeg

wat hij precies voelde
en precies bedoelde

met zijn losbandigheid?
Pennywise ten spijt

zat hij vol kattenkwaad,
maar ook vol goede raad.

Hij veroorzaakt chaos,
maar maakt ook zaken los.

Joker of harlekijn,
Messing vond haar clown fijn.

Zelf manipulatief
had ze haar nar juist lief

om zijn egoïsme
in zijn anarchisme.

Haar manipulatie
leek wel haar protectie

tegen gekwetst worden.
Ongezonde horden

opwerpen tegen zijn
in al haar angst en pijn.

“Liefde zal verdwijnen,
mijn zelf zal wegkwijnen

als ik niet controleer.
Niet meer manipuleer.”

Ze dacht als Anansi,
Hermes of als Loki.

De trickster is kwetsbaar,
maar oogt onaantastbaar.

Handelt uit wantrouwen.
Zal contact afbouwen

als deze te intiem wordt,
trots naar liefdetekort.

Leeft buiten de regels,
het liefst tussen vlegels

en blijft onbereikbaar.
Messing maakte zichtbaar

wat mensen verbergen
door ze steeds te tergen.

De manipulator
schiep haar eigen horror.

Een non had ja-geknikt,
op de vinger getikt

door een boze Messing.
“Graag in je eigen kring

gaan filosoferen.
In dit klooster leren

we over de wijzen
die ons naar God wijzen.

De rest is niet probaat.”
was ooit Messings dictaat.

Na wat gebakkelei
liet de non het er bij.

Ze koos wat doopnamen
die niet vaak voorkwamen,

zodat de verwarring
in de samenleving

over wie men noemde
minder vaak opdoemde.

De abdis was senang
met de frisse ingang

die de non opperde.
Hoewel ze mopperde

stond ze de namen toe.
Ketterij bleef taboe.

De non verbeet zich tam.
De kloosterpriester kwam

toen de doop uitvoeren
bij baby's van boeren.

John was de eerste doop.
Pierre de tweede aanloop.

Gottfried volgde René
en Baruch weer die twee.

De priester zei verheugd:
“Dit doet me als paap deugd.

Vernoemd naar Pierre Bayle.
Gottfried Leibniz ligt me

ook zo na aan het hart.
Net als René Descartes:
“Cogito, ergo sum”.
Zowaar ook John Locke’s roem:

‘De tabula rasa.’
En Baruch Spinoza

zei: “De mens is vrij voor
zover hij alleen door

de rede wordt geleid.”
Het bed is ruim gespreid

voor deze kinderen.
Laat niets hen hinderen.

Bij hun naamgenoten
zijn ze aangesloten!”

Messing had niet geklapt,
toen de non was betrapt.

Wat later zag de non
dat Messing de strijd won.

Het werd een verbanning
wegens een uitspanning.
